\chapter{مبانی نظری و کارهای مرتبط}
\label{chap:background}

\section{ریزشبکه به‌عنوان سامانه تصمیم}
ریزشبکه یک محدوده الکتریکی مشخص است که منابع تولید پراکنده، ذخیره‌ساز و بارها را هماهنگ می‌کند. منابع می‌توانند خورشیدی، بادی، دیزلی یا پیل سوختی باشند. باتری انرژی مازاد را ذخیره و در کمبود، بخشی از بار را تأمین می‌کند. در حالت متصل به شبکه، تبادل با شبکه اصلی هم منبع انعطاف و هم منبع هزینه یا درآمد است.

از دید علوم کامپیوتر، ریزشبکه یک فرایند تصمیم‌گیری ترتیبی است. حالت سامانه دست‌کم شامل وضعیت شارژ و شاید قیمت و پیش‌بینی بار است. کنش‌ها توان‌های قابل کنترل‌اند. پاداش می‌تواند منفیِ هزینه بهره‌برداری باشد. قیود، کنش‌های مجاز را محدود می‌کنند. این نگاه کمک می‌کند بعداً خروجی شبکه عصبی به‌عنوان کنش، و باقیمانده فیزیکی به‌عنوان جریمه، تفسیر شود.

اولادجو و فولی یک ریزشبکه متصل به شبکه را با چند سایت مشارکت‌کننده، تولید خورشیدی و باتری بررسی کرده‌اند. در آن مدل، تبادل انرژی میان مشارکت‌کنندگان و شبکه اصلی در تابع اقتصادی ظاهر می‌شود و مسئله فقط کمینه‌سازی یک هزینه متمرکز نیست؛ توزیع سود نیز اهمیت دارد \pcite{oladejo2019microgrid}. برای پروژه حاضر، دو درس از آن مقاله گرفته می‌شود: فهرست متغیرهای فیزیکی و اقتصادی، و وجود یک لایه اجتماعی-اقتصادی بالای لایه توان.

\section{توازن توان و قرارداد علامت}
ساده‌ترین قانون فیزیکی در لایه مدیریت انرژی، توازن توان در هر گام است. برای مدل پایه می‌توان نوشت:
\begin{equation}
  P_{\mathrm{grid}}(t)+P_{\mathrm{PV}}(t)+P_{\mathrm{dis}}(t)
  -P_{\mathrm{ch}}(t)-P_{\mathrm{load}}(t)=0,
  \label{eq:power-balance}
\end{equation}
که $P_{\mathrm{grid}}$ توان خالص دریافتی از شبکه، $P_{\mathrm{PV}}$ توان خورشیدی، $P_{\mathrm{ch}}$ و $P_{\mathrm{dis}}$ توان شارژ و دشارژ و $P_{\mathrm{load}}$ بار است.

قرارداد علامت باید در کل گزارش و کد یکی باشد. در این متن، مقدار مثبت $P_{\mathrm{grid}}$ یعنی خرید از شبکه و مقدار منفی یعنی فروش به شبکه. شارژ توانی است که به باتری وارد و دشارژ توانی است که از آن خارج می‌شود. اگر این قرارداد در داده برعکس شود، باقیمانده فیزیک درست به نظر می‌رسد اما مدل غلط یاد می‌گیرد. بنابراین یکسان‌سازی واحد و علامت، بخشی از روش است نه کار حاشیه‌ای.

\section{پویایی و حدود باتری}
وضعیت شارژ را می‌توان با رابطه گسسته زیر تقریب زد:
\begin{equation}
  \mathrm{SOC}_{t+1}=\mathrm{SOC}_{t}
  +\frac{\eta_{\mathrm{ch}}P_{\mathrm{ch}}(t)\Delta t}{E_{\max}}
  -\frac{P_{\mathrm{dis}}(t)\Delta t}{\eta_{\mathrm{dis}}E_{\max}},
  \label{eq:soc-dynamics}
\end{equation}
که $E_{\max}$ ظرفیت انرژی، $\Delta t$ گام زمان و $\eta_{\mathrm{ch}}$ و $\eta_{\mathrm{dis}}$ بازده شارژ و دشارژ هستند. حدود بهره‌برداری نیز
\begin{equation}
  \mathrm{SOC}_{\min}\leq \mathrm{SOC}_{t}\leq \mathrm{SOC}_{\max}
  \label{eq:soc-bounds}
\end{equation}
است. معمولاً $\mathrm{SOC}_{\min}$ صفر نیست تا از تخلیه عمیق جلوگیری شود.

مدل‌های دقیق‌تر محدودیت توان مبدل، دمای سلول، پیری و ممنوعیت شارژ و دشارژ هم‌زمان را هم دارند. در نسخه کارشناسی، همین سه عنصر -- توازن، دینامیک گسسته و حدود -- برای ساخت باقیمانده کافی‌اند. اگر دانشجو بخواهد مدل را سنگین‌تر کند، باید همان قید جدید را هم به تابع هزینه یا به تبدیل خروجی اضافه کند.

\section{هدف اقتصادی و مشارکت‌کنندگان}
در لایه اقتصادی، هزینه خرید، درآمد فروش و هزینه استفاده از باتری ظاهر می‌شوند. اگر فقط یک بهره‌بردار وجود داشته باشد، مسئله به کمینه‌سازی هزینه کل نزدیک است. اگر چند ساختمان، فروشگاه یا واحد تولیدی در یک ریزشبکه شریک باشند، کمینه‌سازی جمع کافی نیست؛ باید مشخص شود هر کس چه سهمی از سود می‌برد.

مقاله اولادجو و فولی برای این لایه از راه‌حل چانه‌زنی نش تعمیم‌یافته استفاده می‌کند. ایده این است که سود هر مشارکت‌کننده نسبت به وضعیت عدم‌همکاری بهبود یابد و قدرت چانه‌زنی بتواند نابرابر باشد. برای حل عددی نیز الگوریتم بهینه‌سازی مبتنی بر آموزش و یادگیری (\lr{TLBO}) به کار رفته است که جز اندازه جمعیت و تعداد نسل، پارامتر تنظیمی کمی دارد \pcite{oladejo2019microgrid}.

این لایه برای پروژه فعلی «افق توسعه» است. واردکردن آن در همان تابع هزینه \lr{PINN} بدون داده چندذینفع، ادعا را بزرگ‌تر از شواهد می‌کند. راه درست این است که ابتدا مدل تک‌بهره‌بردار فیزیک‌آگاه پایدار شود و سپس سود ائتلاف با یک قاعده چانه‌زنی روی خروجی‌های可行 تقسیم شود.

\section{شبکه عصبی فیزیک‌آگاه}
رئیسی و همکاران یک معادله دیفرانسیل غیرخطی عمومی را به شکل
\begin{equation}
  u_t+\mathcal{N}[u;\lambda]=0,
  \qquad x\in\Omega,\quad t\in[0,T]
  \label{eq:general-pde}
\end{equation}
در نظر می‌گیرند \pcite{raissi2019pinns}. شبکه $\hat{u}_{\theta}$ پاسخ را تقریب می‌زند. مشتق‌گیری خودکار، مشتق‌های لازم برای ساخت باقیمانده $f$ را می‌دهد. تابع هزینه نوعی آن‌ها ترکیبی از خطای داده، خطای معادله در نقاط هم‌مکان و در صورت نیاز شرایط مرزی است:
\begin{equation}
  \mathcal{L}(\theta)
  =\lambda_{d}\mathcal{L}_{\mathrm{data}}
  +\lambda_{f}\mathcal{L}_{\mathrm{physics}}
  +\lambda_{b}\mathcal{L}_{\mathrm{boundary}}.
  \label{eq:pinn-loss}
\end{equation}

درس قابل انتقال این مقاله، مثال معادله برگرز یا شرودینگر نیست؛ سه اصل است:
\begin{enumerate}
  \item دانش فیزیکی باید به عبارتی قابل ارزیابی روی دسته نمونه تبدیل شود؛
  \item آن عبارت باید نسبت به پارامترهای شبکه مشتق‌پذیر یا دست‌کم برای بهینه‌ساز قابل استفاده باشد؛
  \item وزن و مقیاس مؤلفه‌ها بخشی از طراحی‌اند، نه جزئیات پیاده‌سازی.
\end{enumerate}
در مدیریت انرژی گسسته، بسیاری از قیود اصلاً به مشتق مکانی نیاز ندارند. توازن توان یک باقیمانده جبری است. دینامیک باتری یک باقیمانده یک‌گامی است. این ساده‌سازی نقطه قوت پروژه کارشناسی است، نه ضعف آن.

\section{مشتق‌گیری خودکار و بهینه‌سازی}
کتابخانه‌هایی مانند پای‌تورچ گراف محاسباتی را نگه می‌دارند و با انتشار معکوس، گرادیان تابع هزینه نسبت به وزن‌ها را می‌سازند. اگر ورودی زمان با گزینه نیاز به گرادیان علامت‌گذاری شود، مشتق خروجی نسبت به زمان نیز قابل محاسبه‌ است. برای مدل گسسته این گزارش، آن قابلیت اختیاری است.

رئیسی و همکاران اغلب از ترکیب \lr{Adam} و سپس \lr{L-BFGS} استفاده کرده‌اند \pcite{raissi2019pinns}. همین الگو برای شروع مناسب است: \lr{Adam} برای حرکت اولیه و \lr{L-BFGS} برای تصحیح محلی، در صورتی که اندازه دسته و پایداری عددی اجازه دهد. انتخاب بهینه‌ساز را نباید با موفقیت روش یکی گرفت؛ بدون تنظیم وزن فیزیک، هر دو بهینه‌ساز می‌توانند به پاسخ غیرفیزیکی همگرا شوند.

\section{مقایسه رویکردها}
جدول~\ref{tab:comparison} تفاوت مفهومی مدل داده‌محور، مدل فیزیک‌آگاه و بهینه‌سازی کلاسیک را نشان می‌دهد.

\begin{table}[H]
  \centering
  \caption{مقایسه سه رویکرد برای تصمیم انرژی ریزشبکه}
  \label{tab:comparison}
  \begin{tabular}{>{\centering\arraybackslash}p{0.20\textwidth}>{\centering\arraybackslash}p{0.23\textwidth}>{\centering\arraybackslash}p{0.23\textwidth}>{\centering\arraybackslash}p{0.22\textwidth}}
    \toprule
    \textbf{معیار} & \textbf{شبکه داده‌محور} & \textbf{\lr{PINN}} & \textbf{بهینه‌سازی} \\
    \midrule
    منبع دانش & داده & داده و معادله & مدل صریح \\
    تضمین قید & ضعیف & نرم، با جریمه & معمولاً سخت \\
    نیاز به داده & زیاد & کمتر در ایده & کم، اما مدل دقیق \\
    هزینه آموزش & متوسط & بیشتر & صفر پس از مدل‌سازی \\
    هزینه استنتاج & کم & کم & وابسته به حل مجدد \\
    تفسیر & محدود & مرتبط با فیزیک & شفاف \\
    \bottomrule
  \end{tabular}
\end{table}

هیچ ستونی به‌طور مطلق بهتر نیست. بهینه‌سازی کلاسیک وقتی مدل کامل و زمان حل مجاز باشد قوی است. شبکه داده‌محور وقتی داده فراوان و قید کم‌اهمیت باشد سریع است. فیزیک‌آگاه جایی معنا دارد که داده ناقص است و چند قانون مطمئن در دست است.

\section{جمع‌بندی پیوند دو مقاله}
مقاله \lr{PINN} سازوکار ادغام معادله و داده را می‌دهد \pcite{raissi2019pinns}. مقاله ریزشبکه فهرست متغیر، قید، تبادل با شبکه و لایه تقسیم سود را می‌دهد \pcite{oladejo2019microgrid}. پیوند درست این است که قوانین لایه توان وارد تابع هزینه شوند و قوانین لایه ائتلاف، در صورت وجود چند ذینفع، بعد از تولید برنامه انرژی اعمال شوند. پیوند نادرست این است که اعداد، شکل‌ها یا برتری‌های گزارش‌شده در آن دو مقاله به‌نام این پروژه نوشته شوند.

فصل بعد این پیوند را به تعریف ورودی، خروجی و باقیمانده تبدیل می‌کند.
